\section{Step 10 Move signed values without changing their bits}\label{transpose-s10}

The address equations depend on coordinates, not on the sign of a stored integer. A word load can sign-extend its result in a general register; a word store writes the low 32 bits. These operations preserve a negative input word through the copy. A byte store would preserve only eight bits and leave the rest of the destination wrong.

\subsection{Write or inspect}
Put a negative value and a word with different high and low bytes into the active input. Predict their transposed positions. Compare full words after execution. Restore the starter's index values afterward so the remaining walkthrough uses its stated fixture.

\subsection{Expected observation}
The complete 32-bit patterns appear at the predicted positions, including the sign bit. The input allocation is unchanged.

\Needspace{11\baselineskip}
\subsection{Check your understanding}
Why is a byte-store implementation insufficient?

\begin{itemize}
  \item \textbf{A} It copies only part of each required word.
  \item \textbf{B} It always reverses the bytes.
  \item \textbf{C} It changes the matrix dimensions.
\end{itemize}
